% --- LaTeX Essay Template - S. Venkatraman ---

% --- Set document class and font size ---

\documentclass[letterpaper, 11pt]{article}

% --- Package imports ---

% lipsum is just for generating placeholder text and can be removed
\usepackage{hyperref, lipsum} 

% --- Fonts (Set to Palatino or Times New Roman if desired) ---

% \usepackage{newpxtext, newpxmath}
% \usepackage{times}

% --- Page layout settiings ---

% Set page margins
\usepackage[left=1.2in, right=1.2in, top=.9in, bottom=.9in]{geometry}

% Anchor footnotes to the bottom of the page
\usepackage[bottom]{footmisc}
\usepackage{setspace}
\setstretch{1.15}
% Set line spacing
\renewcommand{\baselinestretch}{1.2}

% Use this for a one-off '*' footnote (e.g., a URL in a heading)
\newcommand{\starfootnote}[1]{%
\begingroup
\renewcommand{\thefootnote}{\fnsymbol{footnote}}%
\footnote{#1}%
\endgroup
\addtocounter{footnote}{-1}%
}

% Set spacing between paragraphs
\setlength{\parskip}{2mm}

% --- Page formatting settings ---



% --- Essay title and author ---
\title{\vspace{-1.5cm} Summary of \textit{Veridical Data Science}, Chapters 12 \footnote{\url{https://vdsbook.com/}}}

\date{\today}

% --- Main text ---
% --- Document starts here ---

\begin{document}
% Set link colors for labeled items and URLs (blue) 
\hypersetup{colorlinks=true, linkcolor=blue, urlcolor=blue}
\maketitle

\section*{Chapter 12: Decision Trees and Random Forests}

Chapter 12 introduces tree models and ensemble learning (Random Forests).

Trees are great data structure which store data in a hierarchical way.
CART, which recursively partition the feature space using greedy optimization algorithms, which select splits that maximize homogeneity of the resulting subsets.
For regression trees, splits are chosen to reduce within-node variation (e.g., variance or squared error); for classification trees, splits are chosen to increase purity using criteria such as Gini impurity or entropy.
Trees are highly interpretable, but single trees can be unstable: small data perturbations or minor preprocessing choices can change the structure dramatically, leading to high variance and poor generalization.
To prevent overfitting, the chapter discusses hyperparameters (depth, minimum leaf size, pruning).

Random forests (RF) is a powerful ensemble method that builds on trees to improve predictability and stability by averaging many trees trained on perturbed versions of the data.
RF introduces randomness through bootstrap resampling of observations for each tree and random subsampling of candidate features.
Since we use bootstrap samples, out-of-bag (OOB) predictions can be used to estimate generalization performance without needing a separate test set.
To interpret RF, we can look at feature importance measures (e.g. permutation importance) to identify which features contribute most to predictions.

Finally, the chapter studies PCS in RF by comparing it with logistics regressions.
Due to the nonlinearity and interactions captured by RF, it can achieve better predictive performance than logistic regression, especially when the true relationship between features and response is complex.
But we still need to be careful about overfitting and interpretability when using RF, and the chapter discusses methods for tuning and interpreting RF models.

% --- Document ends here ---
\end{document}
